\section{Discussion}
\label{sec:discussion}
\paragraph{Inverse Edges}
The addition of inverse edges implies that the agent can also walk from an effect to its cause. In general, their addition has a few benefits, like the possibility to undo wrong actions and to reach nodes that otherwise could not be reached under a given path length constraint. Conversely, they also introduce the possibility to make mistakes through false positives. While these mistakes will always happen for the BFS, our agent can minimize them by pruning the search space, as demonstrated in our experiments in Section~\ref{subsec:evaluation-approach}. As our ablation study in Section~\ref{subsec:ablation-study} shows, adding inverse edges improves performance in practice, so their benefits seem to outweigh the pitfalls. This is in line with prior works that also added inverse edges to improve performance~\cite{Xiong2017DeePpath,Das2018Minerva}.

\paragraph{Open-ended Causal Questions}
Currently, our approach only supports binary causal questions. One idea to answer open-ended questions like \textit{``What causes X?''} or \textit{``What are the effects of X?''} would be the following: At inference time, multiple paths are sampled by the agent.
The entities that occur most often as endpoints of these paths are selected as the answer.

\paragraph{Explainability of Causal Questions}
Our approach has the advantage of not only being able to answer binary causal questions but also providing explanations in the form of paths. The paths indicate mechanisms by which the cause produces the effect, potentially through a chain of multiple entities. Each path might indicate a different mechanism. Moreover, we can use additional provenance data that is part of each relation~\cite{Heindorf2020Causenet}. For example, we can reference the original sentence from which the relation was extracted, which may provide additional insights. Furthermore, each relation contains the URL of the original web source, which can be checked to verify the causal relations and receive further information.

\paragraph{Negative Information} Dealing with negative information both during training and for explanation is an interesting open research challenge. Current causality graphs such as CauseNet do not contain negative edges~\cite{Heindorf2020Causenet} and we train our agent on positive causal question, i.e., questions whose answer is ``yes''.  If the ground-truth answer to a question ``Does X cause Y?'' is ``no'',
and we assume a causality graph with only positive edges,
then it does not matter how our agent walks in the graph, as it will never reach Y.

Regarding explainability, it turns out to be surprisingly difficult to explain why the answer to a causal question is ``no'':
(1)~
One possibility might be KGs with negative edges. However, this only works for 1-hop negative paths as negative edges are not transitive. For example, on a 2-hop negative path like ``A does not cause B does not cause C'', we cannot conclude that ``A does not cause C''.
(2)~
Instead of materializing negative edges, link prediction approaches might be used to predict negative edges on the fly. However, link prediction approaches come with their own challenges, e.g., they rarely make perfect predictions.
(3)~
Complex background knowledge might be encoded in the form of an ontology that allows to express negative statements (e.g., by using negation and forall quantifiers in description logic). However, ontological reasoning comes with its own challenges, e.g., it is hardly scalable to real-world, large-scale knowledge graphs.
